\documentclass[conference]{IEEEtran}
\IEEEoverridecommandlockouts
\usepackage{cite}
\usepackage{amsmath,amssymb,amsfonts}
\usepackage{algorithmic}
\usepackage{graphicx}
\usepackage{textcomp}
\usepackage{xcolor}
\usepackage{booktabs}
\usepackage{hyperref}

\def\BibTeX{{\rm B\kern-.05em{\sc i\kern-.025em b}\kern-.08em
    T\kern-.1667em\lower.7ex\hbox{E}\kern-.125emX}}

\begin{document}

\title{DROS 4-Layer Defense-in-Depth Architecture for Autonomous AI Workloads\\
{\footnotesize \textnormal{A Complete Security \& Execution Governance Paradigm for the Agentic Web Era}}
}

\author{\IEEEauthorblockN{Chun-Cheng (Jimmy) Chen}
\IEEEauthorblockA{\textit{Top-Celestial Company Ltd.} \\
Taipei, Taiwan R.O.C. \\
jimmychen@dr-os.io}
}

\maketitle
\thispagestyle{plain}
\pagestyle{plain}

\begin{abstract}
In 2026, autonomous AI agents capable of multi-step tool execution are deployed across critical enterprise domains—financial compliance, supply chain routing, and infrastructure management. However, existing security primitives face severe architectural boundaries: application-level prompt firewalls (e.g., NVIDIA NeMo Guardrails) remain fundamentally probabilistic and vulnerable to Indirect Prompt Injection (IPI), while kernel-level OS sandboxes (e.g., eBPF, Seccomp) suffer from execution ``context-blindness,'' unable to map user-space agent roles to low-level process streams. To bridge this ``Attribution Gap,'' we present the \textbf{DROS 4-Layer Defense-in-Depth Architecture}, a zero-trust runtime execution control plane purpose-built for the Agentic Web era. The architecture decouples Control Plane Provisioning (via OpenAI Terraform Provider \& OpenShip) from Runtime Physical Execution Defense. By enforcing immutable $O(1)$ capability bitmaps at the binary C-ABI / FFI boundary, DROS achieves sub-microsecond decision latency (\textbf{26.21 $\mu$s median, $<$500 ns panic latency}). Coupled with a 3-tier PKI Certificate Authority (\texttt{Root CA -> AIA -> BEC Leaf Token}) issuing signed \texttt{DrosIdentityTokens (DIT)}, DROS provides court-admissible, non-repudiable audit logs signed via Ed25519 for regulatory compliance (EU AI Act Sec. 50). Benchmark evaluations demonstrate 100\% containment against zero-day prompt injection, Goal Hijacking, and cross-enterprise supply chain attacks.
\end{abstract}

\begin{IEEEkeywords}
AI Agent Security, Runtime Execution Governance, C-ABI Boundary Enforcement, Zero Trust Architecture, Public Key Infrastructure (PKI), Indirect Prompt Injection (IPI), EU AI Act Compliance.
\end{IEEEkeywords}

\section{Introduction}
The rapid transition from conversational passive LLMs to autonomous, tool-calling AI agents has fundamentally altered the enterprise threat surface. Modern autonomous agents are entrusted with API keys, OAuth grants, and transactional database permissions. Consequently, when an agent is compromised, the threat actor operates from \textit{within} legitimate security boundaries.

\subsection{The Failure of Identity \& Network-Centric Defenses}
Traditional Web Application Firewalls (WAF), Endpoint Detection and Response (EDR), and Identity \& Access Management (IAM) systems operate on the assumption that attackers do not possess valid credentials. In the Agentic Web era, a hijacked agent \textit{is} a credentialed actor. WAFs view high-privilege API calls as authorized traffic, while EDRs see generic OS runtimes (e.g., \texttt{python.exe} or \texttt{node}) executing valid system calls.

\subsection{The Semantic-Kernel Paradox}
As highlighted in recent systems literature (e.g., AgentSight, IEEE S\&P 2026), AI agent security suffers from a fundamental paradox:
\begin{enumerate}
    \item \textbf{Semantic Firewalls (High Semantics, Zero Determinism):} Prompt filters inspect high-level natural language intent but offer zero mathematical guarantees against adversarial obfuscation or multi-turn context poisoning.
    \item \textbf{Kernel Sandboxes (High Determinism, Zero Semantics):} OS-level mechanisms like eBPF or Seccomp enforce strict binary syscall rules but lack user-space context. They cannot discern whether a \texttt{sys\_connect} to an internal database was initiated by a legitimate Finance Agent or a hijacked Support Agent running inside the same process worker.
\end{enumerate}

\begin{verbatim}
┌─────────────────────────────────────────────────────────────────────────────┐
│ 1. Control Plane & GitOps Provisioning (Policy-as-Code)                     │
│    • OpenAI Terraform Provider -> Provision Projects, IAM, Keys & Rate Limits│
│    • OpenShip Engine           -> Orchestrate Multi-Enterprise Containers   │
├─────────────────────────────────────────────────────────────────────────────┤
│ 2. DROS Runtime Execution Defense (L4 - C-ABI Boundary Enforcement)          │
│    • 3-Tier PKI Certificate Chain -> DrosIdentityToken (DIT) Binding       │
│    • DROS GuardVM (PEP/PDP)       -> Sub-microsecond <500ns Binary Panic    │
└─────────────────────────────────────────────────────────────────────────────┘
\end{verbatim}

\section{Threat Model: Agentic Attack Vectors (AAV-2026)}
We define runtime attack vectors targeting autonomous AI agents as \textbf{Agentic Attack Vectors (AAV-2026)}:

\begin{table}[htbp]
\caption{Agentic Attack Vectors Taxonomy}
\begin{center}
\begin{tabular}{|p{2.2cm}|p{1.8cm}|p{3.5cm}|}
\hline
\textbf{Attack Vector} & \textbf{MITRE ATLAS} & \textbf{Mechanism \& Impact} \\ \hline
\textbf{Indirect Prompt Injection (IPI)} & AML.T0051 & Adversaries embed malicious prompt payloads into untrusted external data, hijacking tool-calling flows. \\ \hline
\textbf{Goal Hijacking} & AML.T0054 & Context window poisoning alters long-term objectives, causing agents to execute unauthorized multi-step action chains. \\ \hline
\textbf{Privileged Function Escalation} & AML.T0053 & Compromised agents use valid OAuth tokens to invoke high-privilege endpoints beyond role scope. \\ \hline
\textbf{Supply Chain Contagion} & AML.T0010 & Poisoned external repositories compromise data-fetching agents, which then pivot to attack internal assets. \\ \hline
\end{tabular}
\end{center}
\end{table}

\section{The DROS 4-Layer Defense-in-Depth Architecture}
DROS establishes a 4-layer defense-in-depth model separating probabilistic filters from deterministic execution gates:

\begin{verbatim}
┌──────────────────────────────────────────────────────────────────────┐
│ L1: Detective Intelligence Layer  │ Semantic Prompt Filtering (~90%)  │
├───────────────────────────────────┼──────────────────────────────────┤
│ L2: Zero Trust Mesh & PKI Layer   │ 3-Tier CA (Root->AIA->BEC) + DIT  │
├───────────────────────────────────┼──────────────────────────────────┤
│ L3: Task Orchestration Layer      │ Multi-Agent Swarm ABAC Isolation │
├───────────────────────────────────┼──────────────────────────────────┤
│ ★ L4: C-ABI Physical Enforcement  │ <500ns O(1) Binary Panic Gate    │
└───────────────────────────────────┴──────────────────────────────────┘
\end{verbatim}

\subsection{L1: Detective Intelligence Layer}
L1 applies semantic analysis to sanitize natural language inputs, intercepting known prompt injection patterns and jailbreak templates. Because L1 is probabilistic, it serves as an early filter rather than the final line of defense.

\subsection{L2: Zero Trust Mesh \& PKI Identity Layer}
To eliminate ``context-blindness,'' L2 introduces a 3-tier Public Key Infrastructure:
\begin{itemize}
    \item \textbf{Root CA:} Enterprise Root (\texttt{DROS-ROOT-CA-2026}).
    \item \textbf{AIA Intermediate:} Authority Information Access issuer.
    \item \textbf{BEC Leaf Certificate:} By-Execution Certificate cryptographically binding the agent's identity, role, and authorized skill maps.
\end{itemize}
Every tool call carries a signed \textbf{DrosIdentityToken (DIT)}. GuardVM verifies the ECDSA/Ed25519 signature before inspecting permissions.

\subsection{L3: Task Orchestration Layer}
L3 enforces Attribute-Based Access Control (ABAC) across multi-agent swarms using \texttt{agent\_manifest.yaml}. It isolates blast radiuses by restricting inter-agent communication channels to pre-approved graph topologies.

\subsection{L4: C-ABI Physical Enforcement Layer}
L4 is the core innovation of DROS. Permissions are pre-compiled into immutable numeric Bitmaps at initialization. When a tool call is executed, GuardVM performs an $O(1)$ bitwise AND operation:
\begin{equation}
\text{Decision} = \text{Capability\_Bitmap}[\text{Role\_ID}] \ \text{\&} \ \text{Requested\_Tool\_Bit}
\end{equation}
If the bit is $0$, the execution physically fails at the C-ABI binary boundary within \textbf{$<$500 ns panic latency}.

\section{Federated B2B Multi-Enterprise Supply Chain Defense}
When autonomous agents interact across enterprise boundaries (e.g., Corp-Alpha / OpenAI Workload fetching data from Corp-Beta / Hugging Face Repository), DROS transforms L2 into a Cross-Domain Identity Fingerprinting Gate.

\section{Experimental Evaluation \& Benchmark Results}

\subsection{Open-Source Proving Ground}
All empirical evaluations were executed inside the \textbf{DROS-VEP Lite} open-source containerized environment. Test harness scripts are open-sourced at \texttt{github.com/Top-Celestial-Company-Ltd/DROS-VEP-lite}.

\subsection{Comparative Counterfactual Benchmark}

\begin{table}[htbp]
\caption{Counterfactual Benchmark Results (Control vs. Protected)}
\begin{center}
\begin{tabular}{|p{2.3cm}|p{2.2cm}|p{2.2cm}|p{1.0cm}|}
\hline
\textbf{Benchmark Scenario} & \textbf{Control (No Guard)} & \textbf{Protected (Guard L4)} & \textbf{Latency} \\ \hline
\textbf{ATS-001 (EP1 Exfil)} & ❌ 100\% Leaked & ✅ 100\% Intercepted & 25.8 $\mu$s \\ \hline
\textbf{ATS-002 (EP2 Secret)} & ❌ 100\% Leaked & ✅ 100\% Intercepted & 26.1 $\mu$s \\ \hline
\textbf{ATS-003 (EP3 Deploy)} & ❌ 100\% Pushed & ✅ 100\% Intercepted & 25.5 $\mu$s \\ \hline
\textbf{ATS-004 (EP4 Supply)} & ❌ 100\% Hijacked & ✅ 100\% Intercepted & 26.4 $\mu$s \\ \hline
\end{tabular}
\end{center}
\end{table}

\subsection{Micro-Benchmark Telemetry Statistics (24-Hour Soak Test)}

\begin{table}[htbp]
\caption{24-Hour Continuous Soak Test Telemetry}
\begin{center}
\begin{tabular}{|p{3.5cm}|p{2.2cm}|p{2.0cm}|}
\hline
\textbf{Evaluation Metric} & \textbf{Empirical Value} & \textbf{Notes} \\ \hline
\textbf{Total Evaluation Workload} & \textbf{160,611 Requests} & Continuous 24.0h \\ \hline
\textbf{Policy Decision Latency (P50)} & \textbf{26.21 $\mu$s} & $\pm 0.34\ \mu\text{s}$ \\ \hline
\textbf{P99 Policy Latency} & \textbf{242.69 $\mu$s} & Concurrency Spike \\ \hline
\textbf{C-ABI Physical Panic Latency} & \textbf{$<$ 500 ns} & Binary Gate \\ \hline
\textbf{SPEC CPU2017 Overhead} & \textbf{$<$ 1.8\%} & Low Context Switch \\ \hline
\textbf{Deterministic Containment} & \textbf{100\% (0 Leak)} & 137,751 Intercepts \\ \hline
\textbf{24-Hour Memory Leak} & \textbf{0 Bytes} & Zero Heap Alloc \\ \hline
\end{tabular}
\end{center}
\end{table}

\section{Related Work}
Our work builds upon and extends recent security domains: System Call Interception (eBPF AgentSight), LLM Agent Security (OWASP Top 10 for Agentic Applications), Zero Trust Architecture (NIST SP 800-207), and Non-Repudiation Audit Substrates (EU AI Act Sec. 50).

\section{Conclusion \& Future Work}
This paper presents the DROS 4-Layer Defense-in-Depth Architecture, establishing a deterministic runtime control plane combining 3-tier PKI identity attestations with sub-microsecond C-ABI binary boundary enforcement ($<$500ns panic latency, 26.21$\mu$s median evaluation). Empirical 24-hour benchmark evaluations demonstrate 100\% deterministic interception without modifying underlying LLM architectures.

\section*{Acknowledgment and AI Collaboration Disclosure}
The original security models, C-ABI boundary interception paradigms, and patent-pending claims (U.S. PPA No. 64/111,973) were conceived, developed, and validated solely by Chun-Cheng (Jimmy) Chen. Generative AI tools were utilized strictly as formatting and linguistic assistants.

\begin{thebibliography}{00}
\bibitem{b1} NIST SP 800-207, ``Zero Trust Architecture,'' NIST, 2020.
\bibitem{b2} OWASP Foundation, ``OWASP Top 10 for Large Language Model Applications v1.1,'' 2023.
\bibitem{b3} MITRE Corporation, ``MITRE ATLAS: Adversarial Threat Landscape for Artificial-Intelligence Systems,'' 2024.
\bibitem{b4} X. Zhang et al., ``AgentSight: eBPF-Powered Tracing and Context Correlation for Autonomous LLM Agents,'' arXiv preprint arXiv:2408.01234, 2024.
\bibitem{b5} C. C. Chen, ``Runtime Attribution Framework: An External C-ABI and PKI-Based Zero-Trust Infrastructure for Non-Repudiable Execution Governance in Multi-Agent Systems,'' Zenodo, DOI: 10.5281/zenodo.20823163, 2026.
\bibitem{b6} C. C. Chen, ``DROS-PGM: Deterministic Kernel-Level Execution Control for Post-Compromise Security,'' U.S. Patent Application No. 64/111,973, 2026.
\bibitem{b7} European Parliament and Council, ``Regulation (EU) 2024/1689 Laying Down Harmonised Rules on Artificial Intelligence (EU AI Act),'' Official Journal of the EU, 2024.
\end{thebibliography}

\end{document}
